\documentclass[11pt]{article} \usepackage{amsmath,amssymb,amsthm,mathtools} \usepackage{geometry} \usepackage{hyperref} \geometry{margin=1in} \title{Prime-Indexed Self-Consistent Systems:\\ A Unified Framework for Domain Stability, EMIRP Closure,\\ and Invariant Dynamics} \author{Gnome Badhi\\ {\small Independent Systems Analyst}} \date{May 21, 2026} \theoremstyle{definition} \newtheorem{definition}{Definition}[section] \newtheorem{axiom}{Axiom}[section] \theoremstyle{plain} \newtheorem{theorem}[definition]{Theorem} \newtheorem{proposition}[definition]{Proposition} \newtheorem{corollary}[definition]{Corollary} \theoremstyle{remark} \newtheorem{remark}[definition]{Remark} \begin{document} \maketitle \begin{abstract} We introduce the \emph{Prime-Indexed Self-Consistent System} (PISCS), a mathematical framework in which primes index minimal admissible domains and EMIRPs (reversible primes) define representation-invariant validation operators. Each domain $D_p$ carries a high-resolution operator $F_p$, while each EMIRP $e$ induces a validation operator $F_{e^R}$ acting as a projection onto a self-consistent subspace. Stability is achieved when \[ F_{e^R} \circ F_p = F_p, \] so that the high-resolution operator is a fixed point of EMIRP validation. This invariant is expressed equivalently as a commutative diagram, a natural transformation, a dynamical fixed point, a projection kernel, and a renormalization fixed point. The structure provides a minimal mathematical model for self-validation, identity preservation, and hierarchical domain construction. \end{abstract} \section{Introduction} Self-consistent systems require mechanisms for validating their own outputs. Physical theories require scale invariance and closure; cognitive and computational systems require self-monitoring; and mathematical structures require invariants that persist across transformations. In the PISCS framework, primes index minimal admissible domains, and EMIRPs (primes whose digit reversals are also prime) define reversible validation structures. The central idea is: \begin{itemize} \item a prime $p$ defines a stable domain $D_p$; \item an EMIRP $e$ defines a validation operator $F_{e^R}$; \item the EMIRP flip $e \leftrightarrow e^R$ encodes a symmetry that determines a projection geometry; \item stability promotes the domain to a higher-level variable when the invariant \[ F_{e^R} \circ F_p = F_p \] holds. \end{itemize} We develop this structure in several equivalent forms: as a category of admissible domains, as a kernel-based projection geometry, as a dynamical fixed-point condition, and as a renormalization fixed point derived from an action principle. \section{Minimality and Prime-Indexed Domains} \subsection{Admissible domains} Let $P$ denote the set of prime numbers. For each $p \in P$, we associate a domain $D_p$ together with a distinguished operator $F_p : D_p \to D_p$. \begin{definition}[Admissible domain] An \emph{admissible domain} is a pair $(D,F)$ where $D$ is a representation space (e.g.\ a vector space, function space, or configuration space) and $F : D \to D$ is a distinguished endomorphism encoding the intrinsic dynamics or structure of the domain. \end{definition} \subsection{Minimality axiom} \begin{axiom}[Domain minimality] A domain $D_p$ indexed by $p$ is \emph{minimal admissible} if and only if: \begin{enumerate} \item \textbf{Irreducibility:} \[ D_p \not\simeq D_a \times D_b \qquad (a,b>1). \] \item \textbf{Operator closure:} \[ \mathrm{End}(D_p)=\{\mathrm{id}_{D_p},F_p\}, \qquad F_p^2 = F_p. \] \item \textbf{Non-derivability:} \[ D_p \notin \mathrm{Gen}(D_q) \qquad \forall q < p. \] \item \textbf{Unique factorization:} \[ D_n \simeq D_{p_1} \times \cdots \times D_{p_k} \quad\Rightarrow\quad \{p_1,\dots,p_k\} \text{ is unique}. \] \end{enumerate} \end{axiom} \begin{remark} This axiom mirrors the arithmetic role of primes: minimality corresponds to irreducibility, and unique factorization of composite domains parallels the unique factorization of integers into primes. In PISCS, primes thus index the minimal admissible domains $D_p$. \end{remark} \subsection{Prime-indexed domain function} We summarize the association of primes to domains via the function \[ \mathcal{D} : P \to \mathsf{Dom}, \qquad \mathcal{D}(p) = D_p. \] Examples (conceptual, not unique): \[ \mathcal{D}(2)=D_2,\qquad \mathcal{D}(5)=D_5,\qquad \mathcal{D}(7)=D_7. \] Each domain $D_p$ corresponds to a minimal combinatorial or functional structure capable of supporting the operator $F_p$. \section{EMIRP Symmetry and Validation} \subsection{EMIRPs and reversal} Let $\mathcal{E} \subset P$ denote the set of EMIRPs: primes whose digit reversals (in a fixed base, e.g.\ base 10) are also prime. For each $e \in \mathcal{E}$, let $e^R$ denote its reversal. In the abstract PISCS setting, we do not rely on base‑dependent arithmetic; instead, the map \[ e \mapsto e^R \] encodes a structural symmetry that determines a validation operator and a projection geometry. \subsection{Validation operators and projection} To each EMIRP $e \in \mathcal{E}$ we associate a validation operator \[ F_{e^R} : X \to X, \] acting on the representation space $X$ of the relevant domain(s). \begin{axiom}[Validation idempotence] \[ F_{e^R}^2 = F_{e^R}. \] \end{axiom} \begin{axiom}[Validation projection] The operator $F_{e^R}$ is self‑adjoint (with respect to a chosen inner product on $X$) and therefore defines an orthogonal projection onto a closed subspace: \[ S_{e^R} = \mathrm{Im}(F_{e^R}) \subseteq X. \] \end{axiom} \begin{remark} Geometrically, $F_{e^R}$ collapses $X$ onto the \emph{self‑consistent} subspace $S_{e^R}$ determined by the EMIRP symmetry. Vectors in $S_{e^R}$ are fixed points of validation; vectors orthogonal to $S_{e^R}$ are annihilated. \end{remark} \subsection{Promotion and hierarchical variables} Given an EMIRP $e$ and its associated domain $D_e$, we may \emph{promote} the entire domain to a higher‑level variable once it is stabilized by validation. Define the promotion map: \[ \mathcal{V} : \mathcal{E} \to \mathsf{Var}, \qquad \mathcal{V}(e) = v(D_e), \] meaning that the domain $D_e$, once stabilized by $F_{e^R}$, is treated as a single variable at the next level of description. \section{Operator Algebra and the Unified Invariant} Each domain $D_p$ carries a high‑resolution operator \[ F_p : D_p \to D_p, \] and each EMIRP $e$ induces a validation operator \[ F_{e^R} : X \to X. \] The central invariant of the PISCS framework is: \[ F_{e^R} \circ F_p = F_p. \] This expresses: \begin{itemize} \item \textbf{self‑validation} — the domain operator is already consistent with EMIRP symmetry, \item \textbf{representation invariance} — EMIRP reversal does not alter the domain dynamics, \item \textbf{structural closure} — the operator remains stable under validation, \item \textbf{domain stability} — the domain is promoted only when the invariant holds. \end{itemize} Equation $F_{e^R} \circ F_p = F_p$ is the backbone of PISCS; the remainder of the manuscript shows how this invariant appears in multiple mathematical guises. \section{Category-Theoretic Formulation} We now present a categorical formulation of PISCS, showing that the EMIRP validation structure arises naturally as a functor and that the core invariant corresponds to a naturality condition. \subsection{The category of admissible domains} Define a category $\mathcal{C}$ whose objects are the minimal admissible domains: \[ \mathrm{Ob}(\mathcal{C}) = \{ D_p \mid p \text{ prime} \}. \] \begin{definition}[Morphisms] For each prime $p$, the morphisms of $D_p$ are: \[ \mathrm{Hom}(D_p, D_p) = \{ \mathrm{id}_{D_p}, F_p \}, \] with composition rules: \[ F_p \circ F_p = F_p, \qquad F_p \circ \mathrm{id}_{D_p} = F_p, \qquad \mathrm{id}_{D_p} \circ F_p = F_p. \] There are no morphisms between distinct domains: \[ \mathrm{Hom}(D_p, D_q) = \varnothing \qquad (p \neq q). \] \end{definition} Thus each $\mathrm{End}(D_p)$ is a two‑element idempotent monoid. \subsection{The EMIRP functor} Define a functor \[ F : \mathcal{C} \to \mathcal{C} \] by \[ F(D_p) = D_{e^R}, \qquad F(F_p) = F_{e^R}, \] where $e$ is the EMIRP associated with $p$. Functoriality holds because: \[ F(F_p \circ F_p) = F(F_p) = F_{e^R} = F_{e^R} \circ F_{e^R}. \] \subsection{Natural transformation and the invariant} Define a natural transformation \[ \eta : \mathrm{Id}_{\mathcal{C}} \Rightarrow F. \] For each object $D_p$, the component \[ \eta_{D_p} : D_p \to D_{e^R} \] is the canonical embedding of the domain into its EMIRP‑validated counterpart. \begin{theorem}[Naturality condition] For every morphism $F_p : D_p \to D_p$, \[ F(F_p) \circ \eta_{D_p} = \eta_{D_p} \circ F_p. \] \end{theorem} If $\eta_{D_p}$ is chosen to be the identity on the underlying representation space, the naturality condition reduces to: \[ F_{e^R} \circ F_p = F_p, \] which is exactly the PISCS invariant. \subsection{Interpretation} The categorical formulation shows: \begin{itemize} \item EMIRP validation is a functorial lift, \item self‑consistency is a naturality condition, \item the invariant expresses that $F_p$ is a fixed point of the EMIRP functor. \end{itemize} This provides a clean categorical backbone for the entire framework. \section{Kernel–Projection Geometry} We now rewrite and expand the kernel‑operator formulation of PISCS, showing that the EMIRP validation operator acts as a projection onto the self‑consistent subspace. \subsection{Kernel representation of operators} Assume $X$ is a Hilbert space. Each operator admits a kernel representation: \[ (F_p x)(u) = \int K_p(u,v)\, x(v)\, dv, \] \[ (F_{e^R} x)(u) = \int K_{e^R}(u,v)\, x(v)\, dv. \] \subsection{Kernel form of the invariant} The invariant \[ F_{e^R} \circ F_p = F_p \] is equivalent to the kernel identity \[ \int K_{e^R}(u,w)\, K_p(w,v)\, dw = K_p(u,v), \] or simply \[ K_{e^R} K_p = K_p. \] \subsection{Projection geometry} We impose the structural axiom: \[ F_{e^R}^2 = F_{e^R}. \] Together with self‑adjointness, this implies: \begin{proposition} $F_{e^R}$ is the orthogonal projection onto the subspace \[ S = \mathrm{Im}(F_p). \] \end{proposition} Thus: \[ F_{e^R}\big|_{S} = \mathrm{id}_{S}. \] Geometrically: \begin{itemize} \item $F_p$ pushes vectors into $S$, \item $F_{e^R}$ collapses the entire space onto $S$, \item vectors in $S$ are fixed by both operators. \end{itemize} \subsection{Spectral decomposition} Since $F_{e^R}$ is a self‑adjoint projection, its spectrum is: \[ \sigma(F_{e^R}) = \{0,1\}. \] Thus: \[ X = \ker(F_{e^R}) \oplus \mathrm{Im}(F_{e^R}), \] and the kernel has the form \[ K_{e^R}(u,v) = \sum_i \phi_i(u)\, \overline{\phi_i(v)}, \] where $\{\phi_i\}$ is an orthonormal basis of $\mathrm{Im}(F_{e^R})$. \subsection{Fixed‑point subspace} The fixed‑point subspace of $F_p$ is: \[ \mathrm{Fix}(F_p) = \{ x \in X : F_p(x) = x \}. \] If $F_p$ is idempotent, then: \[ \mathrm{Fix}(F_p) = \mathrm{Im}(F_p). \] Thus the fixed‑point subspace is exactly the EMIRP‑stable subspace. \section{Dynamical System Form} We now express the PISCS invariant in dynamical terms. Let $X$ be the representation space of the domain $D_p$, and let \[ x_t \in X \] denote the system state at time $t$. \subsection{High-domain evolution} The high‑resolution operator $F_p$ defines the domain dynamics: \[ x_{t+1} = F_p(x_t). \] This represents the intrinsic evolution of the domain indexed by $p$. \subsection{EMIRP validation dynamics} The EMIRP validation operator $F_{e^R}$ acts as a projection onto the self‑consistent subspace: \[ y_{t+1} = F_{e^R}(x_{t+1}). \] Since $F_{e^R}$ is idempotent and self‑adjoint, $y_{t+1}$ is the closest self‑consistent state to $x_{t+1}$. \subsection{Stability condition} A stable domain satisfies: \[ x_{t+1} = y_{t+1}. \] Thus: \[ F_{e^R}(F_p(x_t)) = F_p(x_t), \] which is exactly the invariant \[ F_{e^R} \circ F_p = F_p. \] \subsection{Interpretation} The dynamical interpretation is: \begin{itemize} \item $F_p$ generates the raw evolution, \item $F_{e^R}$ enforces self‑consistency, \item stability occurs when the raw evolution already lies in the validated subspace. \end{itemize} Thus the invariant expresses that the domain dynamics are already self‑consistent and require no correction. \section{Renormalization Flow} We now express the invariant in terms of renormalization. The renormalization map extracts the scale‑invariant content of the domain operator. \subsection{Renormalization map} Define the renormalization map: \[ \mathcal{R}(F_p) = F_{e^R} \circ F_p. \] This map takes a high‑resolution operator and returns its EMIRP‑validated effective operator. \subsection{Fixed points} A domain is renormalization‑stable when: \[ \mathcal{R}(F_p) = F_p. \] \begin{theorem}[Renormalization fixed point] The operator $F_p$ is a fixed point of the renormalization map if and only if \[ F_{e^R} \circ F_p = F_p. \] \end{theorem} \subsection{Interpretation} The renormalization interpretation is: \begin{itemize} \item $F_p$ is the high‑resolution operator, \item $F_{e^R}$ extracts the scale‑invariant content, \item the invariant states that $F_p$ is already scale‑invariant. \end{itemize} Thus EMIRPs define the projection onto the renormalization fixed‑point manifold. \section{Action Principle} We now derive the invariant from a variational principle. The action formulation shows that the PISCS invariant is not only structural but also variationally optimal. \subsection{Lagrangian} Define the Lagrangian: \[ \mathcal{L}(x) = \big\| F_{e^R}(F_p(x)) - F_p(x) \big\|^2. \] This quantity measures the deviation from self‑consistency: the larger the difference between $F_{e^R}(F_p(x))$ and $F_p(x)$, the more the system violates EMIRP stability. \subsection{Action} Define the action: \[ S[x] = \int \mathcal{L}(x(t))\, dt. \] The action accumulates inconsistency over time. \subsection{Stationarity condition} Stationarity of the action, \[ \delta S = 0, \] requires that the integrand vanish for all admissible variations. Thus: \[ F_{e^R}(F_p(x)) = F_p(x). \] This is precisely the PISCS invariant. \subsection{Interpretation} The action principle interpretation is: \begin{itemize} \item the system minimizes inconsistency, \item the minimum is achieved when the EMIRP projection does nothing, \item the invariant is the Euler–Lagrange condition for extremizing the action. \end{itemize} Thus the invariant is not merely structural — it is the variational optimum of the system. \section{Unified Invariant and Structural Synthesis} We now consolidate the multiple mathematical formulations of the PISCS invariant. Across operator algebra, category theory, kernel geometry, dynamical systems, and renormalization, the same structural condition appears: \[ F_{e^R} \circ F_p = F_p. \] This single equation expresses: \begin{itemize} \item \textbf{self‑validation}: the domain operator is already consistent with EMIRP symmetry, \item \textbf{representation invariance}: EMIRP reversal does not alter the domain dynamics, \item \textbf{projection geometry}: $F_{e^R}$ is the projection onto the fixed‑point subspace of $F_p$, \item \textbf{categorical naturality}: the EMIRP functor commutes with domain morphisms, \item \textbf{renormalization stability}: $F_p$ is a fixed point of the renormalization map, \item \textbf{variational optimality}: the action is minimized when the invariant holds. \end{itemize} The invariant is therefore not an isolated identity but a unifying structural law. \subsection{Equivalence of formulations} \begin{theorem} The following statements are equivalent: \begin{enumerate} \item $F_{e^R} \circ F_p = F_p$. \item $F_{e^R}(F_p(x)) = F_p(x)$ for all $x \in X$. \item $F_{e^R}$ is the orthogonal projection onto $\mathrm{Im}(F_p)$. \item $\eta : \mathrm{Id}_{\mathcal{C}} \Rightarrow F$ is a natural transformation. \item $F_p$ is a fixed point of the renormalization map $\mathcal{R}(F_p)=F_{e^R}\circ F_p$. \item $F_p$ extremizes the action $S[x]$. \end{enumerate} \end{theorem} \begin{proof} (1) $\Rightarrow$ (2): Apply both sides to $x$. (2) $\Rightarrow$ (3): Projection invariance implies $F_{e^R}$ acts as identity on $\mathrm{Im}(F_p)$. (3) $\Rightarrow$ (4): Projection invariance yields naturality. (4) $\Rightarrow$ (5): Functoriality ensures renormalization stability. (5) $\Rightarrow$ (6): Fixed points minimize the action. (6) $\Rightarrow$ (1): Euler–Lagrange equations yield the invariant. \end{proof} Thus the invariant is the structural fixed point of the entire PISCS architecture. \section{Conclusion} The Prime‑Indexed Self‑Consistent System (PISCS) provides a unified mathematical framework in which: \begin{itemize} \item primes index minimal admissible domains, \item EMIRPs define reversible validation symmetries, \item validation operators act as orthogonal projections, \item domain operators are fixed points of EMIRP validation, \item the invariant appears across multiple mathematical lenses. \end{itemize} This framework offers a minimal model for: \begin{itemize} \item identity preservation, \item hierarchical domain construction, \item invariant physical and computational structure, \item self‑consistent dynamics across scales and representations. \end{itemize} The equivalence of categorical, geometric, dynamical, and variational formulations suggests that PISCS captures a fundamental structural principle underlying self‑consistent systems. \appendix \section{Projection Operator Proofs} \subsection{Idempotence and projection} \begin{proposition} If $F_{e^R}$ is self‑adjoint and satisfies $F_{e^R}^2 = F_{e^R}$, then it is the orthogonal projection onto its image. \end{proposition} \begin{proof} Self‑adjoint idempotent operators are orthogonal projections by the spectral theorem. The spectrum is $\{0,1\}$, and the operator acts as identity on the $1$‑eigenspace and annihilates the $0$‑eigenspace. \end{proof} \subsection{Kernel representation} If $F_{e^R}$ is a projection, then its kernel has the form: \[ K_{e^R}(u,v) = \sum_i \phi_i(u)\, \overline{\phi_i(v)}, \] where $\{\phi_i\}$ is an orthonormal basis of $\mathrm{Im}(F_{e^R})$. \section{Naturality and the Invariant} \begin{proposition} The naturality condition \[ F(F_p) \circ \eta_{D_p} = \eta_{D_p} \circ F_p \] is equivalent to the invariant. \end{proposition} \begin{proof} If $\eta_{D_p}$ is the identity on the underlying representation space, the naturality square reduces to: \[ F_{e^R} \circ F_p = F_p. \] \end{proof} \section{Renormalization Fixed Points} \begin{proposition} $F_p$ is a fixed point of the renormalization map \[ \mathcal{R}(F_p) = F_{e^R} \circ F_p \] if and only if the invariant holds. \end{proposition} \begin{proof} Immediate from the definition of $\mathcal{R}$. \end{proof} \section{Action Variation} \begin{proposition} Stationarity of the action \[ S[x] = \int \|F_{e^R}(F_p(x(t))) - F_p(x(t))\|^2\, dt \] implies the invariant. \end{proposition} \begin{proof} The integrand is minimized when the norm vanishes, i.e. \[ F_{e^R}(F_p(x)) = F_p(x). \] \end{proof} \begin{thebibliography}{99} \bibitem{MacLane} S.~Mac Lane, \emph{Categories for the Working Mathematician}, Springer, 1998. \bibitem{ReedSimon} M.~Reed and B.~Simon, \emph{Methods of Modern Mathematical Physics, Vol. I: Functional Analysis}, Academic Press, 1980. \bibitem{Zeidler} E.~Zeidler, \emph{Applied Functional Analysis}, Springer, 1995. \bibitem{PeskinSchroeder} M.~Peskin and D.~Schroeder, \emph{An Introduction to Quantum Field Theory}, Westview Press, 1995. \bibitem{Gelfand} I.~Gelfand and S.~Fomin, \emph{Calculus of Variations}, Dover, 2000. \bibitem{Rudin} W.~Rudin, \emph{Functional Analysis}, McGraw–Hill, 1991. \end{thebibliography} \end{document}
